%% main.tex
%% Documento principal para IEEE Journal
%% Basado en bare_jrnl.tex V1.4b

\documentclass[journal,onecolumn]{IEEEtran}
% *** PAQUETES ***
% Descomentar según necesidad
% \usepackage{cite}
% \usepackage[pdftex]{graphicx}
% \usepackage{amsmath}
% \usepackage{algorithmic}
% \usepackage{array}
% \usepackage[caption=false,font=footnotesize]{subfig}
\usepackage{url}
\usepackage{float}
\usepackage[spanish, es-tabla]{babel}
\usepackage{array}
\usepackage{xcolor}
\usepackage{float}
%\usepackage[table]{xcolor}
\usepackage{graphicx}
\usepackage[version=4]{mhchem}
\usepackage{newunicodechar}
\usepackage[spanish]{babel}
\usepackage{longtable}
\usepackage{colortbl}
\usepackage{hyperref}
\newunicodechar{₂}{$_2$}

%Diagrama de gantt dependencias:
% Gantt chart
\usepackage{pgfgantt}
% Tablas
\usepackage{booktabs}   % \toprule, \midrule, \bottomrule
\usepackage{tabularx}   % entorno tabularx

%% metadata.tex
%% Título, autores y encabezados del documento

% *** TÍTULO ***
\title{Análisis predictivo de riesgo de corrupción en contratación pública mediante aprendizaje automático}

% *** AUTORES ***
\author{
        Aragón~Jhojan,
        Juan Amaya, 
        Santiago Buitrago%
\thanks{Universidad Distrital Francisco Jose De Caldas, Colombia}%
\thanks{ jhsaragonr@udistrital.edu.co}%
\thanks{Documento Marzo 2026}}

% *** ENCABEZADOS ***
\markboth{Metodología de la Investigación,~Universidad Distrital,~Marzo~2026}%
{Aragón \MakeLowercase{\textit{et al.}}: Sistema de Bicicletas Compartidas}

\begin{document}
\maketitle


%% abstract.tex
\begin{abstract}
Este proyecto propone la aplicación de técnicas de aprendizaje automático para el análisis predictivo del riesgo de corrupción en la contratación pública en Colombia. A partir de datos de SECOP II, se seleccionarán variables relevantes para entrenar modelos de clasificación que permitan identificar contratos con alta probabilidad de presentar irregularidades. La metodología incluye validación cruzada y evaluación mediante métricas como AUC-ROC y F1-score. El propósito es desarrollar un sistema de alertas tempranas que fortalezca la vigilancia preventiva y contribuya a una gestión más transparente y eficiente de los recursos públicos.
\end{abstract}

\begin{IEEEkeywords}
Elefantes Blancos, Sistema General de Regalías, Contratación pública, SECOP II, 
Aprendizaje automático, XGBoost, Random Forest, Alertas tempranas, 
Detección de anomalías, Auditoría preventiva
\end{IEEEkeywords}

\IEEEpeerreviewmaketitle

% *** SECCIONES DEL DOCUMENTO ***
%  %% introduction.tex
%% Sección de introducción y tabla de contenido

\tableofcontents

% \begin{enumerate}
% \item Listado de Figuras
% \item Listado de Tablas
% \item Descripción del Documento
% \item Generalidades del Proyecto
% \item Especificación y Recabación de Requerimientos Funcionales
% \item Motivadores Arquitecturales
% \item Contexto
% \item Descripción y Justificación de Patrón(es)
% \item Puntos de Vista y Modelos Arquitecturales
% \item Justificación y Relaciones entre los Puntos de Vista
% \item Firmas de Aceptación con Fecha
% \end{enumerate}

\listoffigures
\listoftables



\section{Descripción del Documento}

*Introduccion*
\section{Identifiación del Problema}
Colombia históricamente ha tenido problemas con la ejecución de proyectos de 
infraestructura que están financiados por el Sistema General de Regalías (SGR) algunos 
de los problemas en los que se ha visto afectada han sido asociados con retrasos, 
sobrecostos, incumplimientos contractuales e incluso el abandono de obras, muchos de 
estos estan relacionados con la corrupción en la que está sumergido este pais, lo que ha 
generado los denominados “Elefantes Blancos” que son todos aquellos proyectos o 
infraestructura costosa de mantener que genera más gastos que beneficios, resultando 
que esta sea inútil y difícil de deshacerse de ella. Esta problemática no solo es generada 
por la ineficiencia en la gestión de los recursos. Sino que también genera impedimentos 
en el desarrollo social y económicos de las regiones beneficiaras. 

Pese a los esfuerzos institucionales para mejorar los mecanismos de control y vigilancia, 
La gran mayoría de las acciones de auditoria continúan siendo de carácter reactivo, es 
decir, toman acción una vez el daño se ha materializado o cuando se ve un contratiempo 
visible en la ejecución del proyecto. Lo anterior reduce la capacidad de las entidades de 
control para la prevención de irregularidades que puede haber en un proyecto público. 

En este entorno, la accesibilidad de grandes volúmenes de datos que provienen de 
plataformas como SECOP II, que es una plataforma transaccional en la cual las 
entidades estatales y los proveedores hacen toda la gestión en línea y en tiempo real  de 
todas las etapas de los procesos de contratación, en ellos encontramos procesos como: 
planeación, selección, ejecución y liquidación, Gracias a estas plataformas se abren 
nuevos caminos para el análisis predictivo a través de técnicas de Machine Learning. 
Pese a lo anterior en la actualidad no se cuenta con alguna herramienta sistemática que 
disponga de variables contractuales y financieras para hacer un predicción o 
anticipación desde la etapa de selección, por lo tanto, hay una alta probabilidad que los 
proyectos presenten contratiempos en su ejecución. 

Debido a esto, surge la necesidad de desarrollar una herramienta que implemente 
modelos predictivos que permita identificar de manera temprana contratos con alto 
riesgo de convertirse en “Elefantes Blancos” y proyectos que puedan estar bajo procesos 
irregulares, contribuyendo a que los entes de control cambien de un enfoque reactivo a 
uno preventivo. Este planteamiento no solo fortalecería la toma de decisiones de las 
entidades de control, sino que también promovería una gestión más eficiente, 
transparente y responsable en el manejo de los recursos públicos en proyectos estatales. 
\section{Justificación}

La contratación pública en Colombia constituye un mecanismo fundamental para la ejecución de proyectos de infraestructura y desarrollo social. Sin embargo, la persistencia de ineficiencias y riesgos en su gestión ha generado impactos negativos tanto en el uso de los recursos públicos como en el bienestar de las comunidades.


Por lo general, los mecanismos de control y vigilancia en la contratación estatal han actuado bajo un enfoque reactivo, actuando una vez el daño patrimonial se ha realizado. Este enfoque limita la capacidad de las entidades de control para prevenir riesgos y optimizar la asignación de recursos en los procesos de auditoría. A raíz de esto, diferentes estudios han planteado la necesidad de cambiar a un esquema de control preventivo que se apoye en el análisis de datos y el uso de tecnologías emergentes, como el machine learning.

En los últimos años, el crecimiento de plataformas de datos abiertos como SECOP II ha permitido el acceso a grandes volúmenes de datos estructurados sobre procesos en la contratación pública. Estos datos incluyen información de todas las etapas en este proceso. La accesibilidad a estos datos ha impulsado el desarrollo de enfoques basados en machine learning para la detección de irregularidades, anomalías relacionadas con corrupción y problemas contractuales.  Investigaciones previas han demostrado que variables como el tipo de contratación, el número de oferentes y el valor del contrato contienen información que puede ser usado para anticipar comportamientos irregulares desde las primeras etapas del proceso de contratación \cite{gallego2021preventing} \cite{salazar2024vigia}.

Asimismo, Varios trabajos han evidenciado la efectividad de modelos de aprendizaje supervisado como Random Forest, regresión logística y XGBoost han facilitado la predicción de riesgos asociados a proyectos de infraestructura, abarcando retrasos y sobrecostos \cite{decarolis2022corruption},\cite{alshboul2022extreme},\cite{galjanic2025neural}. En el entorno colombiano, estudios que incluyen datos de plataformas públicas han identificado factores clave que se relacionan con  desviaciones de tiempo y costo, lo que resalta la viabilidad de utilizar datos abiertos para este tipo de analisis \cite{gomez2025data},\cite{gomez2025analyzing}.



Además, desde una perspectiva institucional, el uso de la inteligencia artificial en la contratación pública ha sido reconocido como una herramienta con gran capacidad para mejorar la transparencia y trazabilidad en las decisiones administrativas de igual forma reduce la irregularidades en la información que favorecen la corrupción \cite{calderon2025contratacion},\cite{sierra2025contratacion}. De este modo, los sistemas de alertas tempranas se plantean como una alternativa que puede complementar a los mecanismos de control actuales, permitiendo centrar los esfuerzos de supervisión en contratos con mayor probabilidad de riesgo, especialmente en aquellos donde los recursos son limitados. \cite{gallego2021preventing},\cite{ospina2024percepcion}.

El presente proyecto se justifica en la necesidad de integrar estos avances teóricos y tecnológicos en una propuesta aplicada al entorno del Sistema General de Regalías (SGR), el cual, concentra una proporción significativa de la inversión pública en infraestructura y ha estado históricamente asociado a problemáticas de ejecución. Pese a la disponibilidad de datos y de los avances en analitica, aun existe una brecha en el desarrollo de herramientas que permitan identificar de manera sistemática, desde la etapa de adjudicación, contratos con alto riesgo de convertirse en proyectos fallidos  .

Con esto, el desarrollo de un modelo predictivo que use técnicas de Machine Learning que sea orientado como un sistema que pueda realizar alertas tempranas, no solo representa un gran avance académico en el area de la analitica de datos aplicados a procesos estatales, sino que también una contribución práctica para fortalecer los mecanismos de control. Este tipo de herramienta puede ayudar a entidades de control y veedurías ciudadanas en la toma de decisiones,  impulsando una gestión más eficiente, transparente y responsable de los recursos públicos,en concordancia con los principios establecidos en el marco normativo de la contratación estatal \cite{ospina2024percepcion}.

Por último, este trabajo se alinea con las tendencias actuales en un gobierno abierto y el uso de datos públicos, donde la analítica avanzada se posiciona como un elemento clave para mejorar la gobernanza y la transparencia en los procesos. Como señalan varios estudios recientes, cruzar datos abiertos con Machine Learning es un avance  para el sector público. Nos permite dejar de ser reactivos y empezar a anticiparnos a los problemas antes de que ocurran \cite{lyra2022fraud}.
\section{Estado del arte}

\subsection{Métodos técnicos basados en Inteligencia Artificial y \textit{Machine Learning}}

Varios investigadores han trabajado en el uso de modelos de aprendizaje automático para detectar problemas en la contratación pública, y sus resultados son bastante relevantes para este trabajo. Por ejemplo, en \cite{gallego2021preventing} crearon un sistema de alertas tempranas usando datos del SECOP que buscaba predecir irregularidades contractuales antes de que ocurrieran, aplicando regresión logística y árboles de decisión. Luego, en \cite{salazar2024vigia} fueron un paso más allá con el sistema VigIA, implementado en la Alcaldía de Bogotá, que usa \textit{Random Forest} e índices de riesgo para detectar contratos con alta probabilidad de presentar demoras o irregularidades financieras. Los dos estudios coinciden en algo importante: las características del contrato que se conocen desde el momento de la adjudicación, como la modalidad, el valor y el número de oferentes, ya permiten anticipar si un contrato va a tener problemas.

También hay trabajos orientados a detectar colusión entre oferentes. En  \cite{tas2024machine} propusieron un algoritmo basado en mínimos cuadrados regularizados por \textit{kernel} (KRLS) que logra identificar colusión en subastas usando solo datos básicos como el precio adjudicado y el número de participantes, sin necesitar información detallada de los pliegos. Por su parte, en \cite{decarolis2022corruption} analizaron indicadores de corrupción en licitaciones de obra vial en Italia, comparando LASSO, Ridge y \textit{Random Forest}, y concluyeron que este último modelo tiene mejor desempeño predictivo, especialmente cuando se incluyen variables que no se reportan de forma estándar.

En cuanto a la predicción de sobrecostos y retrasos en proyectos de infraestructura, también existe bastante literatura. En \cite{gomez2025data} y \cite{gomez2025analyzing} se enfocaron en Colombia en donde aplicaron enfoques basados en datos de plataformas abiertas del Estado para modelar desviaciones de tiempo y costo en proyectos viales y educativos, encontrando que el tipo de proceso de selección, las adiciones contractuales y el origen de los recursos son variables clave. A nivel internacional, podemos ver que en  \cite{alshboul2022extreme} mostraron que \textit{XGBoost} supera a la regresión lineal en la predicción de costos de construcción, gracias a su capacidad para capturar relaciones no lineales. Más recientemente, en \cite{galjanic2025neural} usaron redes neuronales artificiales para predecir retrasos y sobrecostos, logrando un MAPE de 10.93\% para la predicción de retrasos, y destacaron que combinar indicadores cuantitativos y cualitativos mejora notablemente los resultados.

\subsection{Aplicaciones directas en SECOP y contratación pública}

SECOP registra prácticamente todo el ciclo de vida de un contrato público:
modalidad, número de oferentes, valor adjudicado, modificaciones y estado de
ejecución. Esa amplitud lo convierte en la fuente más completa disponible para
estudiar contratación en Colombia, y varios grupos de investigación ya la han
aprovechado. Por ejemplo en \cite{calderon2025contratacion}, \cite{sierra2025contratacion} se agurmeta que trabajar sistemáticamente con esos datos abre la puerta a pasar
del control reactivo (que actúa cuando el daño ya ocurrió) a uno preventivo
basado en análisis predictivo.

El antecedente más directo de este trabajo es el sistema VigIA
\cite{salazar2024vigia}, desarrollado para la Alcaldía de Bogotá. Usando
registros de SECOP~II y modelos de \textit{Random Forest}, construyeron índices
de riesgo que le permiten a la entidad priorizar qué contratos vigilar. El
resultado más útil de ese trabajo no es el modelo en sí, sino lo que demuestra:
la información disponible desde la adjudicación ya contiene señales suficientes
para separar contratos problemáticos de los que no lo son, antes de que empiece
la ejecución.

En una línea similar, en \cite{gomez2025data} analizan proyectos viales en
Colombia y encontraron que tres variables (el tipo de proceso de selección, las
adiciones contractuales y la estructura financiera) explican buena parte de las
desviaciones en tiempo y costo. En \cite{gomez2025analyzing} llegan a conclusiones
parecidas revisando infraestructura educativa. Ambos trabajos usan plataformas
públicas del Estado como fuente de datos, lo que confirma que ese camino es
viable metodológicamente.

Este proyecto sigue esa misma lógica, pero desplaza el foco hacia contratos del
Sistema General de Regalías~(SGR), un universo que concentra montos significativos
y una historia documentada de obras inconclusas o ``elefantes blancos''.


\subsection{Contexto colombiano: corrupción, control institucional y gobierno abierto}

La inteligencia artificial también se ha analizado desde una perspectiva más institucional, pensando en cómo puede apoyar los sistemas de control y transparencia en Colombia. En \cite{calderon2025contratacion} se argumenta que los sistemas predictivos pueden fortalecer el control fiscal preventivo al detectar patrones anómalos en contratos antes de que se conviertan en problemas. Según los autores, esto también mejora la trazabilidad de las decisiones administrativas y reduce las asimetrías de información que favorecen la corrupción.

En esa misma línea, en \cite{gallego2021preventing} se propone que los modelos de alerta temprana son una alternativa complementaria a los enfoques tradicionales de sanción, que actúan después de que el daño ya está hecho. Esto es especialmente relevante en contextos donde los organismos de control no tienen recursos suficientes para auditar todos los contratos.

La apertura de la información contractual genera condiciones que permiten 
a ciudadanos e investigadores monitorear el comportamiento de los actores 
involucrados en la contratación pública \cite{sierra2025contratacion} .

En cuanto a la aceptación institucional, se ha analizado la percepción de 
funcionarios colombianos frente al uso de IA para combatir la corrupción, 
encontrando una valoración positiva de su potencial; sin embargo, persisten 
preocupaciones relacionadas con la transparencia algorítmica y la 
interpretabilidad de los modelos \cite{ospina2024percepcion} . Este aspecto resulta clave al diseñar sistemas de este tipo.

\subsection{Marco normativo y jurídico de la contratación pública}

La contratación estatal en Colombia opera bajo un conjunto de principios
transparencia, economía, responsabilidad y selección objetiva orientados a que
el gasto público sea eficiente y verificable \cite{mendieta2022regimen}. Ese marco
no ha permanecido estático: con el tiempo se fue complementando con plataformas
digitales y políticas de datos abiertos que buscan hacer más trazable cada peso
contratado \cite{sierra2025contratacion,milic2022using}.

Varios autores coinciden en que la transparencia no es solo un valor declarativo
sino una condición operativa para reducir la corrupción y sostener la confianza
ciudadana en las instituciones \cite{calderon2025contratacion,ospina2024percepcion}.Pero en 
\cite{malaret2016reto} se va un paso más allá y plantea que el verdadero reto de la
gobernanza contractual hoy no es aprobar más normas, sino consolidar sistemas de
información abiertos que puedan soportar mecanismos de control más sofisticados.

Desde esa perspectiva, construir un modelo predictivo sobre datos de SECOP~II no
es solo un ejercicio técnico: es coherente con la dirección que el propio marco
normativo señala. La vigilancia preventiva que propone este proyecto es,
en cierto modo, la traducción práctica de los principios que ya están escritos
en la ley.

\subsection{Transparencia y gobernanza de datos abiertos}

Un aspecto que no se puede ignorar al desarrollar sistemas analíticos basados en IA es la calidad y disponibilidad de los datos. En \cite{milic2022using} se propone el indicador \textit{OpenGovB Transparency Indicator}, que sirve para medir qué tan abiertos son los datos gubernamentales de un país, considerando aspectos como accesibilidad, interoperabilidad y reutilización. Este tipo de marcos ayuda a entender qué tan preparada está una institución para aprovechar sus propios datos en iniciativas analíticas.

Además, algunos investigadores han explorado enfoques basados en redes para analizar las relaciones entre actores dentro de los sistemas de contratación. En \cite{wachs2019network} se propone modelar a las empresas participantes como nodos y sus interacciones como aristas en un grafo, lo que permite identificar patrones de comportamiento coordinado que podrían indicar colusión o acuerdos anticompetitivos.


% \input{e)MarcoTeórico}
\input{MarcoTeorico2}
% \section{Objetivos}

\subsection{Objetivo General}
Desarrollar un modelo predictivo basado en técnicas de Machine Learning que funcione como un sistema de alertas tempranas, capaz de identificar contratos públicos de infraestructura con alto riesgo de incumplimiento o asociados a procesos irregulares, en el marco del Sistema General de Regalías (SGR). Para ello, se empleará como fuente principal de información la base de datos de la plataforma SECOP II.

\subsection{Objetivos Específicos}
\begin{enumerate}
    \item Identificar y seleccionar variables relevantes de tipo contractual y financiero —como anticipos, modalidad de contratación, número de participantes y ubicación geográfica— que influyan en el riesgo de incumplimientos e irregularidades en contratos de infraestructura pública.

    \item Recolectar, integrar y preprocesar los datos provenientes de la plataforma SECOP II, garantizando su calidad, consistencia y confiabilidad para el análisis predictivo.

    \item Diseñar y entrenar modelos de aprendizaje supervisado, tales como Random Forest, Regresión Logística y XGBoost, con el fin de predecir de manera precisa factores de riesgo en contratos públicos.

    \item Evaluar el desempeño de los modelos mediante técnicas como la validación cruzada estratificada y el uso de métricas robustas, tales como AUC, precisión, recall y F1-score, para determinar su efectividad.

    \item Comparar los resultados obtenidos por los diferentes modelos con el propósito de identificar aquel que ofrezca la mayor capacidad predictiva en la detección temprana de riesgos.

    \item Proponer un esquema de implementación del modelo como herramienta de alertas tempranas, orientado a fortalecer la detección de anomalías por parte de entidades de control y mecanismos de veeduría ciudadana.
\end{enumerate}
% \section{Alcance del Proyecto}

Este proyecto se centra en construir y evaluar un modelo de aprendizaje automático
capaz de estimar el nivel de riesgo de contratos públicos de infraestructura justo
cuando se adjudican, antes de que comience la ejecución. La materia prima son los
registros abiertos de SECOP~II, particularmente los contratos que se financian con
recursos del Sistema General de Regalías~(SGR).

\subsection{Entregables Funcionales}

Al cierre del proyecto se contará con los siguientes productos concretos:

\begin{itemize}
    \item Un \textit{pipeline} de extracción y preparación de datos que cubra desde
    la descarga de registros de SECOP~II hasta la generación de un conjunto de datos
    listo para modelar. Esto incluye manejo de nulos, corrección de inconsistencias
    territoriales y la construcción indirecta de la variable de riesgo a través de
    indicadores \textit{proxy}.

    \item Un módulo de ingeniería de características con las variables que la
    literatura identifica como más informativas: porcentaje de anticipo, modalidad
    de contratación, número de proponentes y localización del contrato.

    \item Un módulo de entrenamiento que compare tres clasificadores ---Regresión
    Logística, Random Forest y XGBoost--- ajustados con búsqueda aleatoria de
    hiperparámetros (\texttt{RandomizedSearchCV}) para mantenerse dentro de los
    límites computacionales del proyecto.

    \item Un módulo de evaluación con validación cruzada estratificada ($k = 5$) y
    reporte de AUC-ROC, F1-score, precisión y \textit{recall}. Se añadirá un análisis
    de importancia de variables para saber qué factores pesan más en la predicción.

    \item El modelo seleccionado, serializado y documentado, junto con el código
    reproducible del experimento completo.

    \item Un informe técnico que compare los tres algoritmos, discuta qué variables
    resultaron más relevantes y proponga cómo podría escalar este trabajo hacia un
    Sistema de Alertas Tempranas operativo.
\end{itemize}

\subsection{Área de Estudio}

El foco está puesto en contratos de infraestructura del SGR registrados en SECOP~II.
Se eligió ese universo por dos razones prácticas: los montos involucrados son
considerables y la plataforma ofrece datos estructurados de acceso público, lo que
hace viable el proyecto sin necesidad de convenios institucionales.

\subsection{Frontera del Proyecto}

Para no exceder la capacidad del equipo en el tiempo disponible, quedan fuera del
alcance las siguientes actividades:

\begin{itemize}
    \item Desplegar el modelo en producción o integrarlo con los sistemas de alguna
    entidad de control. El trabajo termina en la validación académica del modelo.

    \item Desarrollar cualquier tipo de interfaz de usuario ---web, móvil o de
    escritorio---. Los resultados se presentan en el informe y el repositorio de
    código.

    \item Ejecutar búsquedas exhaustivas de hiperparámetros (Grid Search completo),
    dado el costo computacional que implica sobre XGBoost con espacios de parámetros
    amplios.

    \item Detectar esquemas de corrupción que no dejen huella en los registros
    públicos: acuerdos verbales, testaferros no vinculados o pagos fuera del sistema.
    Eso requiere fuentes de inteligencia financiera que están fuera del alcance de
    este trabajo.
\end{itemize}


% \section{Limitaciones del Proyecto}

Todo modelo predictivo tiene un techo definido por los datos y los recursos con que
se construye. Reconocer esas restricciones desde el inicio no es una debilidad del
proyecto: es una condición para interpretar correctamente sus resultados.

\subsection{Limitaciones Técnicas}

El entrenamiento se hará en Google Colab (GPU compartida) y equipos personales.
Eso impone dos restricciones reales que vale la pena nombrar sin rodeos.

Primero, la memoria disponible por sesión puede ser insuficiente para procesar el
total de registros de SECOP~II de una sola vez. Según el volumen final de datos,
puede ser necesario recurrir a submuestreo estratificado o dividir el procesamiento
en lotes. Ninguna de las dos opciones invalida el experimento, pero sí agrega
un paso de decisión metodológica que habrá que documentar.

Segundo, los espacios de búsqueda de hiperparámetros tendrán que acotarse.
Una búsqueda exhaustiva sobre XGBoost con muchos parámetros puede tomar horas
en los entornos disponibles, lo que hace inviable esa estrategia. La búsqueda
aleatoria (\texttt{RandomizedSearchCV}) es la alternativa estándar en estos casos
y produce resultados competitivos con una fracción del costo computacional, aunque
sin garantía de encontrar la configuración óptima global.

Para que los experimentos sean reproducibles, se fijarán semillas aleatorias y se
registrarán las versiones exactas de todas las librerías en un archivo de entorno.

\subsection{Limitaciones de los Datos}

Las restricciones de datos son más difíciles de mitigar que las técnicas, porque
dependen de cómo está construida la plataforma SECOP~II y de qué información
el Estado colombiano publica de forma abierta.

\subsubsection{No existe una etiqueta oficial de incumplimiento}

SECOP~II no tiene un campo que diga si un contrato terminó mal. Para construir
la variable objetivo hay que inferirla: un contrato se etiqueta como problemático
si acumula varias adiciones, muestra suspensiones reiteradas o si la diferencia
entre la fecha de terminación pactada y la fecha real de liquidación supera un
umbral razonable. Esta aproximación introduce ruido en las etiquetas. Algunos
contratos pueden quedar mal clasificados, lo que afecta la calidad del
entrenamiento supervisado y hace que las métricas del modelo sean probablemente
algo optimistas respecto a un escenario con etiquetas perfectas.

\subsubsection{Calidad irregular de los registros}

Los datos de SECOP~II no son uniformes. Hay campos vacíos, registros duplicados,
municipios codificados de maneras distintas según la entidad que los ingresó y
fechas con formatos inconsistentes. El pipeline de limpieza se encargará de la
mayoría de estos problemas, pero es probable que un porcentaje de contratos deba
descartarse por falta de información mínima. Si ese porcentaje no es aleatorio
---por ejemplo, si las entidades con peor calidad de datos son también las más
propensas a irregularidades--- el conjunto de entrenamiento quedará sesgado.

\subsubsection{El modelo no ve lo que no está publicado}

Al trabajar exclusivamente con datos abiertos, el modelo queda ciego ante
irregularidades que no dejan rastro en SECOP~II. Vínculos societarios entre
contratistas, historial de sanciones en otras entidades o transferencias fuera
del sistema no son capturables con este enfoque. Eso no hace inútil al modelo:
hay suficiente evidencia internacional de que las variables contractuales públicas
sí predicen riesgo. Pero sí fija un límite claro a lo que puede detectar.
% \section{Plan de Trabajo}

El trabajo se organiza en cuatro fases. Las dos primeras preparan los datos y
las variables; las dos últimas construyen, ajustan y comparan los modelos. Algunas
actividades se solapan intencionalmente: el análisis exploratorio arranca en la
fase~2 porque los hallazgos descriptivos orientan qué variables conviene incluir,
y la redacción del informe comienza en la fase~4 antes de que termine el análisis
para no acumular toda la escritura al final.

\begin{table}[H]
  \centering
  \caption{Plan de trabajo por fases y actividades}
  \label{tab:plan}
  \renewcommand{\arraystretch}{1.35}
  \begin{tabularx}{\linewidth}{c l X c}
    \toprule
    \textbf{Fase} & \textbf{Actividad} & \textbf{Descripción / Herramientas} & \textbf{Semanas} \\
    \midrule
    \multirow{2}{*}{1}
      & Recolección y limpieza
      & Descarga de registros SECOP~II vía API o CSV. Normalización de campos,
        imputación de nulos, deduplicación y corrección de codificación territorial.
        Python, Pandas, Requests.
      & 1\,--\,5 \\
      & Variable objetivo
      & Definición del \textit{proxy} de incumplimiento: adiciones contractuales,
        suspensiones reiteradas y desfase entre fecha pactada y fecha real de
        liquidación. Análisis de sensibilidad del umbral elegido.
      & 3\,--\,5 \\
    \midrule
    \multirow{2}{*}{2}
      & Ingeniería de características
      & Selección, transformación y codificación de variables predictoras (anticipos,
        modalidad, proponentes, ubicación). Análisis exploratorio para detectar
        correlaciones y desbalance de clases. Scikit-learn, Feature-engine.
      & 5\,--\,8 \\
      & \textit{Pipeline} de preprocesamiento
      & Construcción del flujo reproducible: imputación, escalado estándar y
        codificación \textit{one-hot}. Serialización con \texttt{joblib} para
        reutilización en inferencia. Scikit-learn Pipeline.
      & 6\,--\,8 \\
    \midrule
    \multirow{2}{*}{3}
      & Entrenamiento de modelos
      & Ajuste de Regresión Logística (línea base), Random Forest y XGBoost.
        Optimización con \texttt{RandomizedSearchCV} ($n\_iter = 50$).
        Google Colab, Scikit-learn, XGBoost.
      & 8\,--\,12 \\
      & Validación y métricas
      & Validación cruzada estratificada $k = 5$. Reporte de AUC-ROC, F1-score,
        precisión, \textit{recall} y curvas PR para cada modelo.
      & 11\,--\,13 \\
    \midrule
    \multirow{2}{*}{4}
      & Análisis comparativo
      & Comparación de modelos. Análisis de importancia con SHAP y
        \textit{permutation importance}. Selección y justificación del modelo final.
      & 13\,--\,15 \\
      & Informe y monografía
      & Redacción del informe técnico y la monografía. Organización del repositorio
        con código, datos de muestra y ambiente reproducible.
      & 14\,--\,16 \\
    \bottomrule
  \end{tabularx}
\end{table}

El repositorio en Git documentará cada experimento con su semilla, versión de
librerías y métricas obtenidas. Eso permite retomar cualquier punto del trabajo
sin tener que repetir el procesamiento completo.
% \section{Cronograma}
\label{sec:cronograma}

\subsection{Diagrama de Gantt}

\begin{figure}[H]
  \centering
  \begin{ganttchart}[
      hgrid,
      vgrid={*{3}{dotted}, black},
      title/.append style={fill=black, draw=black},
      title label font=\color{white}\small\bfseries,
      bar/.append style={fill=gray!40},
      bar label font=\small,
      group/.append style={fill=gray!70, draw=black},
      group label font=\small\bfseries,
      milestone/.append style={fill=black, draw=black},
      x unit=0.55cm,
      y unit title=0.6cm,
      y unit chart=0.55cm,
      bar height=0.5,
      group height=0.3
    ]{1}{16}
    %── Encabezado ─────────────────────────────────────────────
    \gantttitle{Semanas del semestre}{16} \\
    \gantttitlelist{1,...,16}{1} \\
    %── FASE 1 ─────────────────────────────────────────────────
    \ganttgroup{Fase 1: Recolección y Limpieza}{1}{5} \\
    \ganttbar{A1.1\ \ Descarga y consolidación SECOP II}{1}{3} \\
    \ganttbar{A1.2\ \ Limpieza y normalización de datos}{2}{4} \\
    \ganttbar{A1.3\ \ Construcción variable objetivo}{3}{5} \\
    %── FASE 2 ─────────────────────────────────────────────────
    \ganttgroup{Fase 2: Ingeniería de Características}{5}{8} \\
    \ganttbar{A2.1\ \ Selección y construcción de variables}{5}{7} \\
    \ganttbar{A2.2\ \ Pipeline de preprocesamiento}{6}{8} \\
    \ganttbar{A2.3\ \ Análisis exploratorio (EDA)}{6}{8} \\
    %── FASE 3 ─────────────────────────────────────────────────
    \ganttgroup{Fase 3: Entrenamiento y Evaluación}{8}{13} \\
    \ganttbar{A3.1\ \ Regresión Logística (línea base)}{8}{10} \\
    \ganttbar{A3.2\ \ Entrenamiento Random Forest}{9}{11} \\
    \ganttbar{A3.3\ \ Entrenamiento XGBoost}{10}{12} \\
    \ganttbar{A3.4\ \ Validación cruzada estratificada}{11}{13} \\
    %── FASE 4 ─────────────────────────────────────────────────
    \ganttgroup{Fase 4: Análisis y Cierre}{13}{16} \\
    \ganttbar{A4.1\ \ Comparación de métricas y modelos}{13}{14} \\
    \ganttbar{A4.2\ \ Análisis de importancia de variables}{13}{15} \\
    \ganttbar{A4.3\ \ Redacción monografía final}{14}{16} \\
    %── Hito ───────────────────────────────────────────────────
    \ganttmilestone{Entrega final}{16}
  \end{ganttchart}
  \caption{Diagrama de Gantt -- Cronograma del proyecto (16~semanas).}
  \label{fig:gantt}
\end{figure}

\subsection{Tabla de Hitos}

Los hitos funcionan como puntos de control internos. Si al llegar a uno el
entregable no está listo, hay margen para ajustar el ritmo antes de que el
retraso se propague a la fase siguiente.

\begin{table}[H]
  \centering
  \caption{Hitos del proyecto}
  \label{tab:hitos}
  \renewcommand{\arraystretch}{1.3}
  \begin{tabularx}{\linewidth}{c l X c}
    \toprule
    \textbf{N°} & \textbf{Hito} & \textbf{Criterio de cumplimiento} & \textbf{Semana} \\
    \midrule
    H1 & Datos listos
       & Dataset limpio, variable objetivo construida y umbral de \textit{proxy}
         documentado
       & 5 \\
    H2 & Pipeline funcional
       & Flujo de preprocesamiento serializado, probado y versionado en el
         repositorio
       & 8 \\
    H3 & Modelos entrenados
       & Los tres clasificadores ajustados con hiperparámetros registrados
         en el experimento
       & 12 \\
    H4 & Evaluación cerrada
       & Reporte de métricas con validación cruzada completo y revisado
       & 13 \\
    H5 & Modelo seleccionado
       & Elección justificada del algoritmo final con análisis SHAP incluido
       & 15 \\
    H6 & Entrega final
       & Monografía y repositorio de código públicos y documentados
       & 16 \\
    \bottomrule
  \end{tabularx}
\end{table}




\input{conclusion}
% comentado por que daba error: 
%\usepackage{multirow}

% *** AGRADECIMIENTOS ***
\input{Agradecimientos}

% *** REFERENCIAS ***
\ifCLASSOPTIONcaptionsoff
  \newpage
\fi
\nocite{*}
\bibliographystyle{IEEEtran}
\bibliography{IEEEabrv,references}



% *** BIOGRAFÍAS ***
%\input{biographies}

\end{document}